% ============================================================================
% CAPÍTULO 18 — CI/CD\index{CI/CD}: entrega contínua
% Status: texto completo (rodada editorial 2)
% ============================================================================
\chapter{CI/CD: Entrega Contínua}
\label{cap:cicd}

\begin{caixaconceito}
\textbf{CI} (Integração Contínua) roda testes e checagens a cada push;
\textbf{CD} (Entrega Contínua) publica e implanta automaticamente após a
aprovação. Com o \textbf{GitHub Actions}, o pipeline\index{pipeline} vive no próprio
repositório — o padrão de mercado para projetos TypeScript.
\end{caixaconceito}

Você já tem uma suíte de testes (capítulo~\ref{cap:testes}) e uma imagem
Docker (capítulo~\ref{cap:docker}). O problema agora é \emph{garantia}: quem
garante que o código que entrou na \texttt{main} passou nos testes? Quem
garante que o build funciona em uma máquina limpa? Quem implanta? A resposta
deste capítulo: \emph{ninguém decide na hora} — o processo decide, sempre, do
mesmo jeito. É isso que separa projeto de hobby de produto.

Ao final deste capítulo, você será capaz de:
\begin{objetivos}
  \item Escrever workflows do GitHub Actions\index{GitHub Actions} (\texttt{on}, \texttt{jobs}, \texttt{steps});
  \item Rodar lint, typecheck e testes em CI a cada push e Pull Request;
  \item Caching de dependências e artefatos\index{artefatos};
  \item Publicar a imagem Docker no GitHub Container Registry;
  \item Implantar em produção (Vercel/self-host) automaticamente;
  \item Entregar o projeto do capítulo: o pipeline completo do SkillHub.
\end{objetivos}

\section{O fluxo de trabalho moderno}

No fluxo baseado em \emph{trunk} (a branch \texttt{main} é sempre
implantável), cada mudança chega por um Pull Request. E cada \texttt{push} ou
\texttt{pull request} dispara um pipeline com os checks de qualidade:

\begin{enumerate}[leftmargin=*, itemsep=0.4em]
  \item \textbf{Lint} (\texttt{eslint}) + \textbf{formatação} (\texttt{prettier} --check);
  \item \textbf{Typecheck} (\texttt{tsc --noEmit});
  \item \textbf{Testes} unitários e de integração;
  \item \textbf{Build} de produção (pega erros que só aparecem no build);
  \item \textbf{Testes E2E} (Playwright);
  \item \textbf{Publicação} da imagem e \textbf{deploy} (se aprovado).
\end{enumerate}

\begin{caixaconceito}
``CI verde'' é a porta de entrada do Pull Request: nada merge sem passar nos
checks. Isso transforma a qualidade de ``responsabilidade individual'' em
\emph{garantia do processo} — e é exatamente o que as empresas esperam de um
sênior.
\end{caixaconceito}

\begin{caixadica}
Os primeiros cinco checks são \textbf{CI}: respondem ``este código é bom?''
O sexto é \textbf{CD}: responde ``este código chega ao usuário?''. Confundir
os dois é comum — e caro: um pipeline que só testa, mas não entrega, é
integração contínua sem entrega; um pipeline que entrega sem testar é um
acidente esperando acontecer.
\end{caixadica}

\section{Entendendo a sintaxe de um workflow}

Um workflow\index{workflow} do GitHub Actions é um arquivo YAML em
\texttt{.github/workflows/} com três blocos essenciais:

\begin{itemize}[leftmargin=*, itemsep=0.4em]
  \item \texttt{on} — \emph{quando} o workflow roda (push, pull\_request, schedule, workflow\_dispatch);
  \item \texttt{jobs} — \emph{o que} roda, em paralelo por padrão; jobs dependentes usam \texttt{needs};
  \item \texttt{steps} — a sequência de comandos e \emph{actions} (blocos reutilizáveis da comunidade, como \texttt{actions/checkout}).
\end{itemize}

Cada job roda em um runner (\texttt{runs-on: ubuntu-latest} é o padrão), que é
uma máquina virtual limpa provisionada para a execução. É por isso que o CI
pega erros que ``funcionavam na sua máquina'': não há nada instalado além do
que o workflow declara.

\section{Primeiro workflow}

\begin{lstlisting}[style=livro, language=Bash,
  caption={.github/workflows/ci.yml — checagens a cada push/PR.},
  label={list:cicd-ci}]
name: CI

on:
  push:
    branches: [main]
  pull_request:

jobs:
  quality:
    runs-on: ubuntu-latest
    steps:
      - uses: actions/checkout@v4

      - uses: actions/setup-node@v4
        with:
          node-version: 24
          cache: npm

      - run: npm ci
      - run: npm run lint
      - run: npm run check          # tsc --noEmit
      - run: npm test -- --run      # vitest

  e2e:
    runs-on: ubuntu-latest
    services:
      postgres:
        image: postgres:18
        env:
          POSTGRES_USER: app
          POSTGRES_PASSWORD: senha
          POSTGRES_DB: app_test
        ports: ["5432:5432"]
        options: >-
          --health-cmd "pg_isready -U app"
          --health-interval 5s
          --health-timeout 5s
          --health-retries 10
    steps:
      - uses: actions/checkout@v4
      - uses: actions/setup-node@v4
        with:
          node-version: 24
          cache: npm
      - run: npm ci
      - run: npx prisma migrate deploy
      - run: npx playwright install --with-deps
      - run: npm run test:e2e
\end{lstlisting}

\begin{caixadica}
Note o bloco \texttt{services}: o GitHub Actions sobe um PostgreSQL como
\emph{serviço} do job — exatamente como você fez com o Compose no
capítulo~\ref{cap:docker}. O E2E roda contra um banco real, migrado e
saudável, em toda execução. Sem isso, seus testes E2E estariam testando
mocks — e mentindo.
\end{caixadica}

\section{Cache, segredos e ambientes}

\begin{itemize}[leftmargin=*, itemsep=0.4em]
  \item \textbf{Cache}: \texttt{cache: npm} acelera \texttt{npm ci}; Playwright mantém cache próprio dos browsers;
  \item \textbf{Segredos}: \texttt{secrets.DATABASE\_URL} — nunca literais no workflow;
  \item \textbf{Ambientes} (\emph{environments}): \texttt{staging} e \texttt{production} com aprovação manual opcional no deploy\index{deploy};
  \item \textbf{Concorrência}: \texttt{concurrency} evita deploys simultâneos do mesmo branch.
\end{itemize}

\begin{caixaconceito}
\textbf{Segredos} são injetados no runner apenas no momento da execução, e
nunca aparecem no arquivo versionado. O GitHub ainda mascara o valor nos logs.
Se você precisar de um valor sensível no workflow, cadastre-o em
\textbf{Settings $\rightarrow$ Secrets and variables $\rightarrow$ Actions} e
referencie com \texttt{\$ \{\{ secrets\index{secrets}.NOME \}\}} — nunca cole o valor no YAML.
\end{caixaconceito}

\section{Publicação de imagem e deploy}

Depois que qualidade e E2E passam na \texttt{main}, o próximo job publica a
imagem Docker no GitHub Container Registry (GHCR\index{GHCR}) com uma tag imutável — o
SHA do commit:

\begin{lstlisting}[style=livro, language=Bash,
  caption={Job de publicação da imagem no GHCR.}, label={list:cicd-publish}]
  publish:
    needs: [quality, e2e]
    if: github.ref == 'refs/heads/main'
    runs-on: ubuntu-latest
    permissions:
      contents: read
      packages: write
    steps:
      - uses: actions/checkout@v4
      - uses: docker/setup-buildx-action@v3
      - uses: docker/login-action@v3
        with:
          registry: ghcr.io
          username: ${{ github.actor }}
          password: ${{ secrets.GITHUB_TOKEN }}
      - uses: docker/build-push-action@v6
        with:
          context: .
          push: true
          tags: ghcr.io/${{ github.repository }}:${{ github.sha }}
\end{lstlisting}

\begin{itemize}[leftmargin=*, itemsep=0.4em]
  \item \texttt{needs} — este job só roda se os dois anteriores passaram;
  \item \texttt{if} — só na \texttt{main} (não em PRs);
  \item \texttt{permissions} — o mínimo que o job precisa (princípio do menor privilégio);
  \item \texttt{tag por SHA} — cada imagem é única e rastreável até o commit; rollback = apontar para uma tag antiga.
\end{itemize}

\begin{caixadica}
Deploy no Next.js: na Vercel, conectar o repositório já cria o pipeline
(preview por PR + produção em \texttt{main}) — e o GitHub Actions complementa
com qualidade. Para self-host: a imagem Docker acima + um passo de SSH/Helm.
\end{caixadica}

% ============================================================================
\section{Projeto do capítulo: pipeline completo do SkillHub}
\label{sec:projeto18}

\begin{caixaprojeto}
\textbf{Objetivo:} criar o pipeline de qualidade e entrega do SkillHub.

\textbf{Critérios de aceite:}
\begin{itemize}[leftmargin=*, itemsep=0.2em]
  \item Workflow \texttt{ci.yml} com lint, typecheck, testes e build — verde em todo PR;
  \item Job E2E com PostgreSQL como serviço;
  \item Workflow de publicação da imagem no GHCR com tag por SHA;
  \item Deploy automático em staging + produção (Vercel ou self-host) protegido por ambiente;
  \item Badges de status no README do projeto;
  \item Histórico de execuções documentando uma correção de pipeline (bug real corrigido).
\end{itemize}
\end{caixaprojeto}

\subsection{Passo a passo}

\begin{passos}
  \item \textbf{Base}: crie \texttt{.github/workflows/ci.yml} com os checks
  mínimos (lint, typecheck, testes) rodando em todo PR;
  \item \textbf{Build}: adicione um passo que roda o build de produção —
  muitos erros só aparecem no bundle final;
  \item \textbf{E2E}: adicione o job com o serviço PostgreSQL e o Playwright;
  \item \textbf{Branch protection}: no GitHub, torne os checks
  \emph{obrigatórios} para merge na \texttt{main} — o CI só tem poder se
  ninguém conseguir contorná-lo;
  \item \textbf{Publicação}: adicione o job \texttt{publish} (GHCR, tag por SHA);
  \item \textbf{Deploy}: conecte o repositório à Vercel (ou configure o deploy
  self-host) e proteja produção com um \emph{environment};
  \item \textbf{Badges e documentação}: adicione os badges de status ao
  \texttt{README} e registre no histórico a primeira correção real de
  pipeline que você precisou fazer.
\end{passos}

\section{Exercícios de fixação}

\begin{exercicios}
  \item Qual a diferença entre CI e CD?
  \item Liste os 4 checks mínimos de um pipeline de qualidade.
  \item O que faz \texttt{on: pull\_request} no workflow?
  \item Por que segredos nunca devem aparecer literais no workflow?
  \item Escreva um workflow mínimo: checkout, setup-node, npm ci, npm test.
\end{exercicios}

\section{Desafios}

\begin{desafios}
  \item \textbf{Deploy condicional}: publique a imagem só quando a branch for \texttt{main} e o commit tiver tag semver.
  \item \textbf{Matrix}: rode os testes em Node 22 e 24 (matrix) e explique o valor de testar múltiplas versões.
  \item \textbf{Notificação}: adicione um passo que posta o resultado do pipeline no Discord/Slack via webhook.
\end{desafios}

\section{Exercícios de nível sênior}

\begin{seniores}
  \item \textbf{Artefatos e diffs}: gere o bundle de produção como artefato e publique o diff de tamanho a cada PR (guarda contra regressão de performance).
  \item \textbf{Deploy com rollback}: configure deploy atômico com rollback automático em caso de health check falho.
  \item \textbf{Tempo do pipeline}: reduza o tempo total do pipeline em 50\% (cache, paralelismo, jobs separados) e documente o antes/depois.
  \item \textbf{Supply chain}: use \texttt{dependabot} + \texttt{sbom} (CycloneDX) para atender ao \texttt{A03} do OWASP (falhas na cadeia de suprimentos).
\end{seniores}

\section{Armadilhas comuns}

\begin{itemize}[leftmargin=*, itemsep=0.4em]
  \item Testes que só rodam em CI e falham na máquina local (e vice-versa);
  \item Segredos vazados em logs de pipeline;
  \item Deploy direto na \texttt{main} sem ambiente de staging;
  \item Pipeline de 30 min com cache mal configurado;
  \item \texttt{latest} sem tag imutável — impossível voltar atrás;
  \item Check obrigatório configurado depois do deploy (a proteção vem primeiro);
  \item Workflows gigantes em um único job em vez de jobs paralelos independentes.
\end{itemize}

\section{Resumo do capítulo}

\begin{itemize}[leftmargin=*, itemsep=0.3em]
  \item CI roda qualidade a cada push; CD publica e implanta;
  \item GitHub Actions: \texttt{on}/\texttt{jobs}/\texttt{steps} com serviços (banco) e cache;
  \item Segredos via \texttt{secrets}; ambientes protegem produção;
  \item Imagens publicadas com tag imutável por SHA;
  \item Projeto entregue: pipeline completo e documentado.
\end{itemize}

\section{Referências oficiais para aprofundar}

\begin{itemize}[leftmargin=*, itemsep=0.3em]
  \item GitHub Actions Documentation: \url{https://docs.github.com/actions}
  \item Referência de workflow: \url{https://docs.github.com/actions/reference/workflow-syntax-for-github-actions}
  \item GitHub Container Registry: \url{https://docs.github.com/packages/working-with-a-github-packages-registry/working-with-the-container-registry}
  \item GitHub — Dependabot: \url{https://docs.github.com/code-security/dependabot}
\end{itemize}
